\chapter{関連研究}
\chapoverview{本章では，本研究に関連する先行研究を整理する．まず，特許データを用いた技術分析に関する研究を概観し，特許情報が技術動向やイノベーション分析に利用されてきたことを確認する．次に，特許文書埋め込み，技術機会分析，クロスドメイン技術探索に関する研究を整理する．最後に，本研究が，全文類似ではなく課題文脈と解決手段文脈を分けて扱う点に位置づくことを示す．}

\section{特許データを用いた技術分析}
特許は，発明の公開情報であると同時に，技術変化を観察するためのデータソースでもある．Grilichesは，特許統計がイノベーション指標として利用できる一方，産業や企業によって特許化傾向が異なるため解釈に注意が必要であることを指摘した\cite{griliches1990patent}．Trajtenbergは，特許引用数が発明価値を反映しうることを示した\cite{trajtenberg1990penny}．Jaffeらは，特許引用情報を用いて知識スピルオーバーの地理的局在を分析した\cite{jaffe1993geographic}．

これらの研究は，特許データが技術分析に有用であることを示した．一方で，特許件数や引用数だけでは，どの課題に対してどの解決手段が対応しているのかを十分に説明できない．本研究は，特許を単なる件数や引用数ではなく，課題文脈と解決手段文脈を含むテキストデータとして扱う点で，これらの研究を補完する．

\section{特許文書埋め込みに関する研究}
自然言語処理技術の発展により，特許文書をベクトル空間上に写像し，意味的類似性を分析する研究が進んでいる．Sentence-BERTは，文単位または文書単位の意味的類似度を効率的に計算する手法として提案された\cite{reimers2019sentencebert}．PatentSBERTaは，特許文書に適した文書埋め込みモデルとして，特許間の距離計算，検索，分類に応用されている\cite{bekamiri2021patentsberta}．

特許文書埋め込みは，類似特許検索や特許分類に有用である．しかし，全文ベクトルのみを用いる場合，課題，解決手段，対象物，効果が混在する．そのため，クロスドメイン技術候補を抽出する際に，候補抽出理由の説明が難しくなる．本研究では，特許文書を全文ベクトルとして扱うだけでなく，課題文脈と解決手段文脈に分離して扱う．

\section{技術機会分析と類推型特許マイニング}
技術機会分析では，特許情報を用いて新たな技術開発の可能性を探索する．類推型特許マイニングでは，ある分野の課題に対して，別分野で使われている解決手段を参考にすることが試みられてきた\cite{jeong2014analogy}．また，BERTとTRIZを組み合わせ，技術要素を拡張しながら技術機会を探索する研究も行われている\cite{wang2023berttriz}．

これらの研究は，異分野の知識を活用して技術機会を探索する点で本研究と関連する．一方で，本研究は，参照分野で高出現である一方，対象分野では低出現である技術手法に注目し，過去期間と将来期間を分けた後ろ向き評価を行う点に特徴がある．

\section{クロスドメイン技術探索}
クロスドメイン技術探索では，ある分野で用いられている技術や知識を，別分野の課題解決に応用する可能性を検討する．Cross-cutting patent analysisでは，対象技術と参照技術の関係を構築し，技術機会からアイデア生成へつなげることが試みられている\cite{liu2021crosscutting}．

本研究では，クロスドメイン技術候補を，単なる類似特許ではなく，課題文脈が近い分野間において，対象分野では低出現である解決手段として定義する．これにより，候補を「なぜ出したのか」を説明しやすくする．

\section{本研究の位置づけ}
\tabref{tab:related_position}に，関連研究と本研究の差分を示す．

\begin{table}[htbp]
  \centering
  \caption{関連研究と本研究の位置づけ}
  \label{tab:related_position}
  \begin{tabularx}{\textwidth}{p{32mm}X X}
    \toprule
    関連領域 & 主な内容 & 本研究との差分 \\
    \midrule
    特許統計・引用分析 & 特許件数や引用数を用いて技術動向や知識スピルオーバーを分析する & 件数や引用数だけでなく，課題文脈と解決手段文脈を含むテキスト情報を扱う \\
    特許文書埋め込み & BERTやPatentSBERTaにより特許間の意味的類似性を分析する & 全文類似だけでなく，課題文脈と解決手段文脈を分ける \\
    技術機会分析 & 特許情報から新たな技術開発機会を探索する & 対象分野で低出現の他分野技術手法を候補化する \\
    類推型特許マイニング & 類似課題を持つ他分野特許を探索する & 課題類似だけでなく，過去期間と将来期間を用いた後ろ向き評価を行う \\
    クロスドメイン技術探索 & 異分野間の技術接続を探す & 課題類似度を前提に，専門家確認候補として提示する \\
    \bottomrule
  \end{tabularx}
\end{table}

本研究の新規性は，BERTを用いること自体ではない．本研究の新規性は，特許文書を課題文脈と解決手段文脈に分離し，課題が近い異分野間において，対象分野では低出現である技術手法を専門家確認候補として抽出する点にある．
